\section{Goal Classifiers：从目标示例学习奖励}

\subsection{基本思想}

收集成功状态（正例）和失败状态（负例），训练一个二分类器，用其输出作为奖励信号。

\begin{enumerate}
    \item 收集成功状态集 $D^+$ 和失败状态集 $D^-$
    \item 训练二分类器：输入 $s_i$，标签 $\mathbf{1}(s_i \in G)$
    \item 用分类器输出作为 RL 的奖励
\end{enumerate}

\begin{warningbox}{Reward Hacking}
RL 算法会积极寻找使分类器输出高分的状态。它可能找到分类器\textbf{从未训练过的}状态——即利用分类器的弱点，而不是真正完成任务。
\end{warningbox}

\subsection{对抗训练解决 Reward Hacking}

解决思路：将 RL 访问的状态加入负例集合，迫使分类器不断更新以避免被利用。

\begin{enumerate}
    \item 初始化 $D^+$（成功状态）和 $D^-$（失败状态）
    \item 用 $D^+$ 和 $D^-$ 训练分类器（平衡数据集）
    \item 用策略 $\pi$ 收集经验 $s_t, a_t, \ldots$
    \item 用分类器奖励更新策略
    \item 将访问的状态加入负例：$D^- \leftarrow D^- \cup \{s_t\}$
    \item 重复
\end{enumerate}

\begin{knowledgebox}{为什么这个方法能工作}
分类器不再能被利用。关键问题：如果策略真的成功了怎么办？只要保持 $D^+$ 和 $D^-$ 的平衡采样，分类器对真正的成功状态仍会输出 $p \geq 0.5$。
\end{knowledgebox}

实验表明，在机器人任务中，基于 50 个演示训练的 goal classifier + RL 达到了 62\% 的成功率，远高于直接模仿学习的 26\%。

\subsection{本章小结}

Goal classifier 是一种简单有效的奖励学习方法。对抗训练是防止 reward hacking 的关键。

